\begin{lemma}[The genuinely outside-prefix reconstruction branch]
Fix the graph, palette, distinguished edge, input, exact terminal tree,
generating coloring, ETT, base, counterexample \(C\), minimum context,
and replacement presentation \(P\). Let \(S\) be a
\noderef{P2880_T1_SplitSwapData} for \(\alpha,\beta,\sigma\).
Write \(\mu=P.\mathit{coloring}\), \(H=P.\mathit{counterexample.hierarchy}\),
\(q=H.\mathit{dividerCount}\), \(A=\noderef{HierarchyStarredStage}(H,q)\),
\(T_{n,q}=H.\mathit{stage}(q)\), and let \(T_n,D_n\) be the terminal
ladder tree and computed Tashkinov stored set of \(P\)'s replacement ETT.
Assume that neither ordered pair \((a,b)=(\alpha,\beta)\) nor
\((a,b)=(\beta,\alpha)\) satisfies
\[
 (q>0\ \text{and }a\in\overline\mu(T_{n,q}))
 \quad\text{or}\quad
 (q=0\ \text{and }[a\in\overline\mu(T_n)\ \text{or }
 (a\in\overline\mu(A),\ a\notin D_n,\ b\notin D_n)]).
\]
There is a \noderef{CorrespondingExtendedTashkinovTree} on the full
\(T\) under \(\sigma\), relative to \(P\)'s replacement ETT, with
maximum property, a hierarchy base, and a good hierarchy having
\noderef{SameGoodHierarchy} with \(P.\mathit{counterexample.good}\).
\end{lemma}
\begin{proof}
Put \(\mathcal E=P.\mathit{corresponding.replacementETT}\),
\(g=P.\mathit{counterexample.good}\), and
\(X=\noderef{HierarchyFullTailPrefix}(C,j-1)\).
Every reference to old maximum property or induction below uses
\(P.\mathit{maximum}\) and \(P.\mathit{minimumContext}\), whose
types are supplied by \noderef{HierarchyReplacementPresentation}.
The \(S\) fields give \(A\preccurlyeq X\preccurlyeq T\) and
\(|A|<|X|\). Apply \noderef{GoodHierarchyTerminalPrefixRepackage}
to \(\mathcal E\), its base and good hierarchy, and these exact
prefix relations. It supplies a maximum ETT \(\mathcal E_X\) on
\(X\), the same rung and certified series, the same core, closure,
and numerical support position, and a good hierarchy \(H_X,g_X\).
It identifies every retained ordinary and starred stage, stored set,
stage-missing set, Gamma pair, and divider position with its old value.
The terminal-kind equality follows from its series equality, rather
than from a claimed equality of only hierarchy stages.

The singleton root set \(\{\operatorname{inl}(y_\ell)\}\), the
swap equality of \(S\), and its exclusion of every vertex-node of
\(X\) form an \noderef{OutsideKempeRecoloring} on \(X\).
Apply \noderef{HierarchyLEMStable} to the certified
\(\mathcal E_X,H_X,g_X\), not to the full \(T\).
Positivity follows from rung equality and the old counterexample.
For its lower-rung premise, rewrite the strict rung inequality by
\(\mathcal E_X.\mathit{rounds}=\mathcal E.\mathit{rounds}\)
and apply the full lower-rung field of the old minimum context.
Properness and exact-one-uncoloredness are the corresponding context
fields. The conclusion supplies
\(\noderef{StableColoring}(X,D_n,\mu,\sigma)\), invariance on \(X\),
and a corresponding maximum prefix ETT with its good hierarchy.
In particular the stability used in this argument is obtained on the
certified prefix; the component is not asserted to avoid full \(T\).

By \noderef{HierarchyBase}, the terminal ladder lies in the core:
in PE this is the left segment of its \noderef{OrderedTreeUnion};
in SE/RE the core is the ladder itself. The core lies in \(X\)
by \(S\). Every incident edge of \(T_n\) is thus incident with
\(X\), and its missing vertices lie there. Restrict the just obtained
stable relation to \(T_n\): the protected set \(D_n\) is unchanged,
and any edge controlled on \(T_n\) is controlled on \(X\), since
tree-missing colors are monotone under vertex-set inclusion.
Rewrite the prefix series equality to identify its \(D_n\) with
that of \(\mathcal E\). Also rewrite
\(\mathcal E.\mathit{generatingColoring}\) to identify its terminal
series coloring with \(\mu\). We now have the precise terminal
stability premise for \noderef{Td} on the original full intermediate
\(\mathcal E\). That theorem, with its positive rung, maximum
property, complete old lower-rung package, and \(S\)'s proper and
one-uncolored fields, supplies \noderef{ModColoring}(\mathcal E,\sigma).

Concatenate the \(S\) prefix route to \(X\) with its new suffix
\(X\to T\) in the reflexive-transitive relation
\noderef{TashkinovAugmentationSequence}. In PE the source is \(T_n\).
In SE/RE retain exactly the supplied edge \(f_n\), singleton
connecting set, and \noderef{TreeSequenceAdjoinEdge} seed. The full
route starts from that seed, even at \(j=2,q=0\) when it equals
\(X\); no nonexistent route from \(T_n\) is substituted.

Suppose the actual terminal kind is PE. Both \(T_n\) and the base
closure have all their vertices in the core, by the two vertex-set
containment clauses of \noderef{OrderedTreeUnion}. The incident-edge
fixation and missing-set equality of \(S\) on \(X\) therefore
give \noderef{StableColoring} on the ladder with protected set \(D_n\)
and on the closure with the empty protected set. Together these are precisely
\noderef{CombinedStableColoring} for \(\mathcal E\)'s original base.
Apply \noderef{HierarchySplitter} to that full \(\mathcal E\),
its base and hierarchy, old maximum and positivity, the old lower-rung
package, the properness and one-uncoloredness of both colorings, this
combined stability, and the full PE route just constructed.
Keep the corresponding ETT, maximum witness, base, and hierarchy from
this single application together. Its conclusion gives
\noderef{SameHierarchyStages} with \(H\), so in particular its
core and PE closure are the old ones. We never identify this certificate
with the independently returned certificate on \(X\), nor with an
arbitrary certificate selected from mod-coloring.

In the actual SE/RE alternative of the \(S\) route, apply the universal
clause of \noderef{ModColoring} to full \(T\).
The strict ladder-to-\(T\) inclusion is the original
\noderef{ExtendedTashkinovTree} field, and the eligibility premise is
the exact seed and concatenated route above. Select a full corresponding
certificate \(c\) relative to \(\mathcal E\).
\noderef{MaximumPropertyTransport} gives maximum property for
\(c.\mathit{replacementETT}\).
Rung, series, ladder, and kind agreement in
\noderef{CorrespondingExtendedTashkinovTree} identify its terminal
ladder with \(T_n\) and retain its non-PE kind.
Construct its hierarchy base with core equal to \(T_n\), closure equal
to the old closure, and support position transported across the exact
ladder equality. Support positivity is retained, the terminal-missing
field is reflexive because its generating coloring is \(\sigma\),
the parallel implication has a false premise, and the nonparallel
equation is exactly the chosen core equality.
Copy \(H\)'s divider count, every exact stage (including terminal
\(T\)), and every numerical divider position. The ladder equality
proves the base-stage field; all nested and strict relations, divider
relations and terminal equality are the unchanged old ones; the core
equality proves core-inside-first. These fields construct a hierarchy
with \noderef{SameHierarchyStages} to \(H\).

We verify goodness for the coherent full hierarchy selected in either
branch. Its core, all retained stages, and in PE its closure have the old
vertices and entries. All retained stages lie in \(A\) or \(T_n\),
and hence in \(X\). Missing sets on these vertices are unchanged.
\noderef{P2752_T2_TerminalStoredColors} identifies the computed terminal
Tashkinov stored set of the full corresponding certificate with \(D_n\).
Consequently \noderef{HierarchyStoredColorsAt}, which subtracts the
starred-stage missing set from \(D_n\), is unchanged at every retained
index. Copy each old Gamma pair, with indices transported by the divider
count equality. Cardinality two and pairwise disjointness follow from
the old \noderef{GoodHierarchy} fields on these same stored sets.

For the missing-color part of avoidance, expand
\noderef{HierarchyStageMissingColors}: each of its vertex sets is
\(T_n\) or a retained ordinary stage inside \(X\), so its missing
union is unchanged. For internal segment avoidance, the next stage
lies inside \(X\); its exact edge colors are fixed. Its truncated
first-missing cutoff is unchanged by
\noderef{HierarchyFirstMissingPositionPrefixCutoff} and equality of
the missing predicate on that exact stage. Thus the internal
\noderef{HierarchySegmentColors} set is identical.
For the terminal segment, fix \(\eta\in D\) and a color \(c\) of
its copied Gamma pair. The first position of \(\eta\) on full
\(T\) is unchanged by \(S\). If \(c\) is not an affected color,
an edge colored \(c\) after the swap already had color \(c\)
before it, directly from \noderef{TwoColorComponentSwap}.
If \(c\) is affected, \(S\)'s Gamma-cutoff field puts its owner's
cutoff at most \(|X|\); every edge before that cutoff is fixed by
\(S\). In both cases membership in the new terminal segment would
imply membership in the old segment, contradicting old avoidance.

For unique refresh, the same stored sets and the same vertex difference
between successive retained stages, with unchanged missing sets, make
the old unique owner satisfy the new predicate and make any new owner
an old one. All equalities of unrefreshed Gamma pairs remain identical.
For each internal minus-closed field, every edge incident with its stage
is fixed by \(S\); its boundary color-edge sets and missing sets are
unchanged, as is the exception union. Hence the old minus-closed
equation applies color by color.
Finally, if PE, every edge incident with the base core is fixed and its
missing colors are unchanged, so the old failure of minus-closedness
with its unchanged Gamma exception union is retained. The ladder and
closure missing sets are unchanged, so the old four-color intersection
bound in the parallel-base field is retained as well. The terminal-kind
agreement is essential here: the new PE premise is exactly the old PE
premise, not an inference from stage equality.
In SE/RE that parallel-base implication is vacuous.
These arguments verify every \noderef{GoodHierarchy} field.
The exact stages, core, PE closure, stored sets, and copied Gamma pairs
then give \noderef{SameGoodHierarchy}, completing the asserted tuple.
\end{proof}
